\chapter{基于Chebyshev多项式近似的多尺度图状态空间模型}

前述章节提出的GGAM和ASD-SSM主要在空域进行特征学习，通过邻域聚合和序列建模捕获几何语义信息，取得了显著性能提升。然而，这些方法的感受野相对固定，难以在同一模块中同时建模不同空间尺度的几何信息。

基于图拉普拉斯的多项式建模为解决上述问题提供了新的视角。通过Chebyshev多项式递推，可以将点云特征分解为多个空间尺度的分量：低阶分量感受野小，捕获点自身及直接邻域的局部几何差异；高阶分量感受野更大，聚合更广范围的上下文结构信息。这种多尺度分解带来三方面优势：首先，能够针对不同空间尺度采用差异化建模策略，避免统一处理导致的建模冲突；其次，各尺度分量的语义纯净度更高，避免了局部细节与全局结构混合建模导致的相互干扰；第三，与空域方法具有天然互补性，两者结合能够更全面地刻画点云的多尺度特征。

为此，本章提出基于Chebyshev多项式近似的多尺度图状态空间模型（Chebyshev Multi-Scale Graph State Space Model, Chebyshev多尺度SSM），与前述空域方法形成互补。Chebyshev多尺度SSM的核心思想是：利用Chebyshev多项式对图拉普拉斯算子进行递推展开，将点云特征分解为多个空间尺度的分量，然后使用独立的Mamba模块分别对各分量进行序列建模，最后通过可学习的加权融合得到增强后的特征表示。这种"多尺度分解 + 分量独立建模"的设计，使模型能够针对不同空间尺度采用差异化的建模策略，从而更有效地捕获点云的多尺度几何信息。

Chebyshev多尺度SSM具有以下优势：首先，通过Chebyshev多项式递推计算，避免了计算密集的特征分解，将复杂度从$O(N^3)$降低至$O(KN)$（其中$K$为多项式阶数，通常为2-5）。其次，通过为每个空间尺度分量配备独立的Mamba模块，使模型能够学习适应不同尺度特性的建模策略。第三，通过可学习的融合权重，模型能够自适应地平衡不同尺度分量的贡献。后续将围绕方法设计和实验分析进行详细介绍。

\section{方法设计}

Chebyshev多尺度SSM的整体流程包括四个核心模块：（1）序列化窗口图构建，利用空间填充曲线的局部性高效构建邻域图；（2）Chebyshev多项式多尺度展开，将点云特征展开为多个空间尺度的分量；（3）分量独立建模，使用独立的Mamba处理各尺度分量；（4）多尺度融合，通过可学习权重自适应融合多尺度特征。整体架构如\cref{fig:chebyshev_arch}所示，接下来将依次介绍各模块的设计细节。

\begin{figure}[htbp]
	\centering
	\includegraphics[width=\linewidth]{pinyu.jpg}
	\caption{Chebyshev多尺度SSM架构图}
	\label{fig:chebyshev_arch}
\end{figure}

\subsection{序列化窗口图构建}

多尺度分解需要首先构建点云的邻域关系。如上文所述，基于序列化的邻域构建方法可以将复杂度从$O(N\log N)$降至$O(N)$。与GGAM采用的窗口划分策略不同，本章采用序列邻居策略，更适合图多项式递推所需的稀疏连接。

具体地，我们利用Point Transformer V3已有的空间序列化顺序，直接使用序列邻域构建图：
\begin{equation}
\mathcal{E} = \{(i, j) \mid q_j \in \mathcal{N}(q_i)\}
\end{equation}
其中$\mathcal{N}(q_i) = \{q_{i-P/2}, \ldots, q_{i-1}, q_{i+1}, \ldots, q_{i+P/2}\}$为点$q_i$在序列上前后各$P/2$个点。本章沿用$P$表示局部邻域大小。对于序列边界点，采用截断策略处理邻域不足的情况，即边界点的实际邻居数可能少于$P$。边权重采用高斯核：
\begin{equation}
w_{ij} = \exp\left(-\frac{\|\mathbf{q}_i - \mathbf{q}_j\|^2}{2\sigma^2}\right)
\end{equation}
其中$\sigma$为尺度参数，设为当前批次所有边长度的均值。该权重设计使得空间距离近的点具有更强的连接权重，距离远的点影响快速衰减，为后续多尺度分解提供合理的邻域拓扑结构。

对于包含多个场景的批次数据，我们为每个场景独立构建图，避免跨场景的信息泄漏。该方法的复杂度为$O(PN) = O(N)$（$P$为常数），且由于序列化顺序已预先计算，无需额外排序开销。

\subsection{Chebyshev多项式多尺度分解}

获得邻域边集$\mathcal{E}$后，下一步是对点云特征进行多尺度分解。设输入特征为$X \in \mathbb{R}^{N \times C}$，我们使用Chebyshev多项式计算$K$个尺度分量的特征表示。

基于边权重$w_{ij}$，构造邻接权重矩阵$\mathbf{W} \in \mathbb{R}^{N \times N}$，其中$\mathbf{W}_{ij} = w_{ij}$若$(i,j) \in \mathcal{E}$，否则为0。度矩阵$D \in \mathbb{R}^{N \times N}$为对角矩阵，对角元素$D_{ii} = \sum_{q_j \in \mathcal{N}(q_i)} w_{ij}$表示点$q_i$的加权度。

为适配Chebyshev多项式对输入特征值范围$[-1, 1]$的要求，需要对归一化拉普拉斯$L_{norm} = I - D^{-1/2}\mathbf{W}D^{-1/2}$进行线性变换。由于$L_{norm}$的特征值范围为$[0, 2]$，构造$\tilde{L} = L_{norm} - I$可将特征值映射到$[-1, 1]$，其中$I$为单位矩阵。代入$L_{norm}$的定义，得到：
\begin{equation}
\tilde{L} = L_{norm} - I = -D^{-1/2}\mathbf{W}D^{-1/2}
\end{equation}

对于输入特征矩阵$X \in \mathbb{R}^{N \times C}$，我们递推计算$K$个Chebyshev基作用的结果：
\begin{equation}
\begin{aligned}
T_0(\tilde{L})X &= X \\
T_1(\tilde{L})X &= \tilde{L}X \\
T_\ell(\tilde{L})X &= 2\tilde{L}T_{\ell-1}(\tilde{L})X - T_{\ell-2}(\tilde{L})X, \quad \ell = 2, \ldots, K-1
\end{aligned}
\end{equation}

关键计算是$\tilde{L}X$，即归一化拉普拉斯与特征矩阵的乘法。利用$\tilde{L} = -D^{-1/2}\mathbf{W}D^{-1/2}$的稀疏性，该操作可高效实现：
\begin{equation}
(\tilde{L}X)_i = -\sum_{q_j \in \mathcal{N}(q_i)} \frac{w_{ij}}{\sqrt{D_{ii}D_{jj}}} X_j
\end{equation}

由于使用scatter操作向量化计算，其复杂度为线性。不同阶的Chebyshev多项式对应不同的空间感受野：$T_0(\tilde{L})X = X$保留点的原始特征（0跳，感受野最小）；$T_1(\tilde{L})X = \tilde{L}X$聚合1跳邻域内的加权差异信息；$T_\ell(\tilde{L})X$的感受野随阶数$\ell$增大而扩展，高阶分量能够捕获更大范围的结构信息。这种多尺度分解为后续的差异化建模奠定了基础。

\subsection{多尺度分量独立建模}

获得$K$个空间尺度的分量$\{T_\ell(\tilde{L})X\}_{\ell=0}^{K-1}$后，本文为每个分量配备独立的Mamba模块，实现尺度特定的序列建模。这一设计源于不同空间尺度分量的固有特性差异：对于低阶分量（$\ell=0,1$），感受野小，捕获点自身及直接邻域的局部几何信息，局部差异和边界细节更为关键；而对于高阶分量（$\ell \geq 2$），感受野大，聚合更广范围的上下文信息，长程依赖关系更为重要。

传统方法使用单一Mamba处理所有尺度的混合特征，难以同时适应小感受野分量的局部特性和大感受野分量的全局特性。本文通过为每个尺度分量配备独立的Mamba，使模型能够学习特定尺度的最优建模策略，从而更有效地捕获不同空间尺度的几何特征。

对于第$\ell$个分量的特征$X^{(\ell)} = T_\ell(\tilde{L})X \in \mathbb{R}^{N \times C}$，沿用Point Transformer V3已有的序列化顺序$Q = (q_1, q_2, \ldots, q_N)$，将其按该顺序排列为序列后送入Mamba进行序列建模。记$X^{(\ell)}_{q_i} \in \mathbb{R}^{C}$为点$q_i$在第$\ell$个分量中的特征向量，则序列化后的输入为：
\begin{equation}
X^{(\ell)}_{seq} = \left(X^{(\ell)}_{q_1},\, X^{(\ell)}_{q_2},\, \ldots,\, X^{(\ell)}_{q_N}\right)
\end{equation}
所有分量使用相同的序列化顺序，保证了分量间的空间对应关系。

第$\ell$个分量的Mamba模块定义为：
\begin{equation}
Y^{(\ell)} = \text{Mamba}_\ell(X^{(\ell)}_{seq})
\end{equation}
其中$\text{Mamba}_\ell$为第$\ell$个分量专属的Mamba块，包含双向SSM（前向+后向）和残差连接。在参数设计上，各分量Mamba的参数完全独立，$\text{Mamba}_\ell$与$\text{Mamba}_j$（$\ell \neq j$）不共享权重。每个分量输出后应用LayerNorm，稳定训练过程。Mamba输出按序列化的逆序恢复为原始点云顺序，得到$Y^{(\ell)} \in \mathbb{R}^{N \times C}$。

\subsection{多尺度特征融合}

获得$K$个尺度分量的Mamba输出$\{Y^{(\ell)}\}_{\ell=0}^{K-1}$后，需要将它们融合为统一的特征表示。本文采用可学习加权融合策略，使模型能够自适应地平衡各尺度分量的贡献。

定义可学习的融合权重$\boldsymbol{\beta} = [\beta_0, \beta_1, \ldots, \beta_{K-1}] \in \mathbb{R}^K$，初始化为均匀权重$\beta_\ell = 1/K$。融合特征计算为：
\begin{equation}
Y_{multi} = \sum_{\ell=0}^{K-1} \alpha_\ell Y^{(\ell)}, \quad \alpha_\ell = \frac{\exp(\beta_\ell)}{\sum_{r=0}^{K-1}\exp(\beta_r)}
\end{equation}
其中$\alpha_\ell$为归一化后的权重（通过softmax保证$\sum_\ell \alpha_\ell = 1$）。这种软加权融合允许梯度反传到所有分量，避免了硬选择的不可微问题。

最终输出通过残差连接和MLP层得到：
\begin{equation}
\begin{aligned}
Z &= \text{Concat}(X, Y_{multi}) \\
\hat{X} &= \text{MLP}(Z) \\
X' &= X + \hat{X}
\end{aligned}
\end{equation}
其中$\text{Concat}(X, Y_{multi}) \in \mathbb{R}^{N \times 2C}$将原始特征和多尺度特征拼接，$\text{MLP}$为2层感知机（隐藏层维度为$C$，使用GELU激活），输出维度为$C$。在MLP的各层之间加入LayerNorm以稳定训练，并使用Dropout防止过拟合。最后通过残差连接得到增强后的特征$X' \in \mathbb{R}^{N \times C}$。

该设计有两个优势：（1）拼接操作保留了原始特征信息，避免多项式变换导致的信息损失；（2）残差连接使模型能够自适应地选择多尺度增强的强度，在不同编码层级和数据集上具有更好的适应性。

与ASD-SSM类似，Chebyshev多尺度SSM也采用层级自适应配置。在编码器-解码器网络的不同层级，多尺度建模的需求不同。对于浅层（第1-2层），特征的空间分辨率高，点云的局部几何信息丰富，启用Chebyshev多尺度SSM带来显著增益。而对于深层（第3-5层），特征已高度抽象，多尺度分解的边际收益递减，可选择性关闭以节省计算。

实验中，我们在前3个层级启用Chebyshev多尺度SSM，后2个层级使用标准Mamba，取得了性能与效率的最佳平衡。

\section{消融实验}

通过系统消融实验，分析Chebyshev多尺度SSM中关键设计选择对模型性能的影响，重点考察多项式阶数、Mamba独立性以及与空域方法的互补性等因素。所有实验均在S3DIS数据集上进行。除"与空域方法的互补性"实验外，其余消融实验均在仅启用Chebyshev多尺度SSM、不加载GGAM和ASD-SSM的条件下进行，以隔离各设计选择的独立影响。

\subsection{Chebyshev多项式阶数}

Chebyshev多项式阶数$K$决定了多尺度分解的精细度。表\ref{tab:cheb_order}展示了不同$K$值的性能对比。

\begin{table}[htbp!]
\centering
\caption{Chebyshev多项式阶数消融实验}
\label{tab:cheb_order}
\begin{tabular}{@{}cccc@{}}
\toprule
阶数$K$ & 参数量 & 推理时间 (ms) & mIoU (\%) \\
\midrule
1 & 1.00$\times$ & 103.9 & 70.5 \\
2 & 1.76$\times$ & 140.9 & 70.7 \\
\textbf{3} & \textbf{2.52$\times$} & \textbf{201.8} & \textbf{70.8} \\
4 & 3.28$\times$ & 250.8 & 70.8 \\
5 & 4.04$\times$ & 306.9 & 71.0 \\
\bottomrule
\end{tabular}
\end{table}

实验结果表明，$K=1$（仅原始特征+1跳邻域）性能最低（70.5\%），说明单一空间尺度不足以捕获复杂的几何结构。从$K=1$到$K=3$，性能随阶数增加而稳步提升（70.5\%→70.8\%），验证了多尺度分解的有效性。$K=4$时性能与$K=3$持平，$K=5$虽略有提升（71.0\%），但参数量增至4.04倍、推理时间增至306.9 ms，相对$K=3$的额外开销远大于0.2\%的边际增益。综合性能与效率，$K=3$为最佳平衡点。

\subsection{Mamba独立性}

本文的核心创新是为每个尺度分量配备独立的Mamba。表\ref{tab:mamba_sharing}对比了不同参数共享策略的效果。

\begin{table}[htbp!]
\centering
\caption{Mamba参数共享策略对比}
\label{tab:mamba_sharing}
\begin{tabular}{@{}lcc@{}}
\toprule
参数共享策略 & 参数量 & mIoU (\%) \\
\midrule
所有分量共享Mamba & 1.00$\times$ & 70.4 \\
低阶/高阶各自共享（2组） & 1.59$\times$ & 70.7 \\
\textbf{每个分量独立Mamba} & \textbf{2.18$\times$} & \textbf{70.8} \\
\bottomrule
\end{tabular}
\end{table}

实验显示，三种策略的推理时间几乎相同（约200 ms），说明计算瓶颈在于图构建和Chebyshev基计算，Mamba本身并非性能瓶颈。因此，独立Mamba仅增加参数量而不增加推理开销。所有分量共享单一Mamba时，参数量最少但难以同时适应不同空间尺度分量的特性差异。将低阶分量（$\ell=0,1$）和高阶分量（$\ell=2$）分为两组各自共享Mamba，性能有所提升，验证了分组建模的价值。每个分量使用独立Mamba时达到最优性能，虽然参数量增至2.18倍，但由于不影响推理速度，这一开销是可接受的。

\subsection{与空域方法的互补性}

表\ref{tab:comparison_spatial}对比了Chebyshev多尺度SSM与空域方法的性能，验证两类方法的互补关系。

\begin{table}[htbp!]
\centering
\caption{Chebyshev多尺度SSM与空域方法对比}
\label{tab:comparison_spatial}
\begin{tabular}{@{}lcc@{}}
\toprule
方法 & 建模策略 & mIoU (\%) \\
\midrule
Baseline（纯Mamba基线） & - & 70.3 \\
+ GGAM & 空域 & 72.5 \\
+ GGAM + ASD-SSM & 空域 & 75.2 \\
+ Chebyshev多尺度SSM & 图多项式 & 70.8 \\
\midrule
\textbf{+ GGAM + ASD-SSM + Chebyshev多尺度SSM} & \textbf{空域+图多项式} & \textbf{75.9} \\
\bottomrule
\end{tabular}
\end{table}

实验表明，单独使用Chebyshev多尺度SSM相比Baseline提升0.5\%（70.3\%→70.8\%），幅度有限，说明图多项式建模单独作用时受限于序列化信息损失等因素。然而，将其与空域方法（GGAM + ASD-SSM）结合后，性能从75.2\%提升至75.9\%（+0.7\%），展现出两类方法的互补性：空域方法通过邻域聚合和多尺度序列建模捕获局部几何细节，Chebyshev多尺度SSM通过图拉普拉斯递推聚合不同跳数的邻域信息，提供了空域方法未覆盖的结构特征，两者结合能够更全面地刻画点云的多尺度几何特征。


\section{本章小结}

本章针对现有空域方法感受野相对固定、难以在同一模块中同时建模不同空间尺度几何信息的局限，提出了基于Chebyshev多项式近似的多尺度图状态空间模型（Chebyshev多尺度SSM）。该方法利用Chebyshev多项式递推计算实现$O(KN)$复杂度的多尺度分解，将点云特征展开为$K$个不同感受野的分量，通过为每个分量配备独立Mamba模块实现尺度特定的序列建模，并通过可学习权重自适应融合多尺度特征。

实验结果表明，Chebyshev多尺度SSM在S3DIS数据集上达到70.8\% mIoU（相比Baseline提升0.5\%），消融实验验证了分量独立建模的必要性。与空域方法（GGAM + ASD-SSM）结合后，完整的PointSS系统在S3DIS上达到75.9\% mIoU，验证了图多项式建模与空域建模的互补性。
